% 来源：1020-12.pdf，第 1--6 页。沿用照片的 hbar=1 约定。
\section{角动量算符}
\subsection{角动量算符的定义}
本章沿用照片中的 $\hbar=1$ 约定。若保留单位，则本章的 $\hat{\bm J}$ 对应物理角动量除以 $\hbar$，其本征值与阶梯系数均为无量纲量。

角动量算符 $\hat{\bm J}=\left(\hat J_x,\hat J_y,\hat J_z\right)$ 是三个厄米算符组成的矢量算符，定义式为
\begin{equation}
\boxed{\hat{\bm J}\times\hat{\bm J}=\mathrm{i}\hat{\bm J}.}
\end{equation}
这里叉乘按固定算符次序定义；由于各分量不对易，它不必为零。分量式为
\begin{equation}
\left[\hat J_x,\hat J_y\right]=\mathrm{i}\hat J_z,\qquad
\left[\hat J_y,\hat J_z\right]=\mathrm{i}\hat J_x,\qquad
\left[\hat J_z,\hat J_x\right]=\mathrm{i}\hat J_y.
\end{equation}
即 $x,y,z$ 循环。用指标表示为
\begin{equation}
\boxed{\left[\hat J_i,\hat J_j\right]=\mathrm{i}\sum_k\epsilon_{ijk}\hat J_k,}
\qquad i,j,k=x,y,z,
\label{eq:angular-commutator}
\end{equation}
其中 Levi--Civita 符号为
\begin{equation}
\epsilon_{ijk}=\begin{cases}
1,&ijk=xyz,yzx,zxy,\\
-1,&ijk=xzy,yxz,zyx,\\
0,&\text{其他情况，即至少两个指标相同。}
\end{cases}
\end{equation}

\begin{example}[单粒子系统中的角动量]
典型例子包括轨道角动量 $\hat{\bm L}=\hat{\bm r}\times\hat{\bm p}$、自旋角动量 $\hat{\bm S}$，以及总角动量 $\hat{\bm J}=\hat{\bm L}+\hat{\bm S}$。在本章单位约定下，它们均满足式\eqref{eq:angular-commutator}；总角动量的结论还利用了轨道与自旋自由度彼此对易。
\end{example}

\subsection{角动量平方算符}
定义
\begin{equation}
\boxed{\hat J^2=\hat{\bm J}\cdot\hat{\bm J}
=\hat J_x^2+\hat J_y^2+\hat J_z^2.}
\end{equation}
可证
\begin{equation}
\boxed{\left[\hat J^2,\hat J_i\right]=0,\qquad i=x,y,z.}
\label{eq:angular-casimir-commutes}
\end{equation}
例如利用 $\left[AB,C\right]=A\left[B,C\right]+\left[A,C\right]B$，有
\begin{align}
\left[\hat J^2,\hat J_z\right]
&=\left[\hat J_x^2,\hat J_z\right]+\left[\hat J_y^2,\hat J_z\right]\notag\\
&=-\mathrm{i}\left(\hat J_x\hat J_y+\hat J_y\hat J_x\right)
+\mathrm{i}\left(\hat J_y\hat J_x+\hat J_x\hat J_y\right)=0.
\end{align}
其余两个分量由循环置换得到。

\subsection{上升算符和下降算符}
定义
\begin{equation}
\boxed{\hat J_+=\hat J_x+\mathrm{i}\hat J_y,\qquad
\hat J_-=\hat J_x-\mathrm{i}\hat J_y.}
\end{equation}
$\hat J_+$、$\hat J_-$ 一般都不是厄米算符，而是互为厄米共轭：
\begin{equation}
\hat J_+^\dag=\hat J_-,\qquad\hat J_-^\dag=\hat J_+.
\end{equation}
由式\eqref{eq:angular-casimir-commutes}，有
\begin{equation}
\left[\hat J^2,\hat J_\pm\right]=0.
\end{equation}
乘积 $\hat J_+\hat J_-$ 和 $\hat J_-\hat J_+$ 都是厄米算符。三个算符的对易关系为
\begin{equation}
\boxed{\left[\hat J_z,\hat J_+\right]=\hat J_+,\qquad
\left[\hat J_z,\hat J_-\right]=-\hat J_-,\qquad
\left[\hat J_+,\hat J_-\right]=2\hat J_z.}
\label{eq:angular-ladder-commutators}
\end{equation}
逆变换为
\begin{equation}
\hat J_x=\frac{\hat J_++\hat J_-}{2},\qquad
\hat J_y=\frac{\hat J_+-\hat J_-}{2\mathrm{i}}.
\end{equation}
因此
\begin{align}
\hat J^2&=\frac12\left(\hat J_+\hat J_-+\hat J_-\hat J_+\right)+\hat J_z^2\notag\\
&=\hat J_+\hat J_--\hat J_z+\hat J_z^2
=\hat J_-\hat J_++\hat J_z+\hat J_z^2.
\end{align}
整理得到后面常用的恒等式
\begin{equation}
\boxed{\hat J_+\hat J_-=\hat J^2-\hat J_z\left(\hat J_z-1\right),\qquad
\hat J_-\hat J_+=\hat J^2-\hat J_z\left(\hat J_z+1\right).}
\label{eq:angular-ladder-products}
\end{equation}
